% ======================================================================
% Chapter 1 -- Introduction
% Target: ~1 page (single section, no subsections).
% PST "classic drama" arc: setting -> problem -> why existing means fall
% short -> approach -> how it solves it (+ contributions) -> outline.
% Present tense, active voice, "we"/"one" (not "I"). VLM not LLM,
% "frozen" not "fine-tuning".
% ======================================================================

\chapter{Introduction}
\label{chap:intro}

Vision-language models acquire broad visual and semantic priors from
large-scale pretraining, yet they are optimised to produce text rather
than to act. Converting that perceptual competence into the ability to
pursue goals inside an interactive environment remains an open problem.

A common route turns a vision-language model into an agent by
post-training the entire backbone on action data, as JARVIS-VLA does for
Minecraft \cite{li2025jarvisvla}. This ties the model's knowledge to a
single embodiment and demands substantial compute. Full fine-tuning is
also risky in itself: updating a strong pretrained backbone can distort
the very features it depends on and weaken out-of-distribution behaviour
\cite{kumar2022finetuningdistortpretrainedfeatures}. Established Minecraft
agents show the two prevailing alternatives, either training a large
control policy from scratch on human video \cite{baker2022vpt} or
bypassing pixel control through symbolic scripting interfaces. Neither
isolates how much agent competence already resides in a frozen,
general-purpose backbone.

This thesis introduces R1V-A, an agent that keeps the vision-language
backbone entirely frozen and trains only a small action head, on the
order of five million parameters, by behavioural cloning on recorded
Minecraft gameplay (rendered by a STEVE-1/VPT agent, not a human; see
Chapter~\ref{chap:method}). Each decision is conditioned on a single rendered
frame together with a natural-language task prompt, and no reward signal
is used during training.

Because the backbone never changes, its features can be computed once and
reused, which makes training cheap enough to sweep several configurations.
The frozen design also turns the backbone into a controlled variable,
which we exploit to make three contributions. First, we compare a joint
vision-language backbone (LLaVA-1.5) against a contrastive one (CLIP) used
purely as frozen feature extractors. Second, we ablate the language
pathway by training each backbone with and without the task prompt, asking
whether the prompt adds signal over image-only control. Third, we evaluate
the resulting agents not only on held-out frames but in live MineRL
rollouts, which probe the out-of-distribution behaviour that frame-level
metrics overlook.

Concretely, the thesis poses three research questions, one per contribution:
\begin{itemize}
\item \textbf{RQ1 (backbone).} Used only as a frozen feature extractor for a
small behavioural-cloning head, does a joint vision-language backbone
(LLaVA-1.5) yield a more capable Minecraft agent than a contrastive one (CLIP)?
\item \textbf{RQ2 (language).} Does conditioning the head on a natural-language
task prompt measurably change the agent's behaviour, and does the answer depend
on the backbone?
\item \textbf{RQ3 (evaluation).} Do live in-environment rollouts reveal
differences between these agents that held-out frame-level metrics do not?
\end{itemize}

The remainder of this thesis is organised as follows.
Chapter~\ref{chap:background} reviews the imitation-learning setting the
agent operates in, Minecraft as a control benchmark, the frozen
vision-language backbone and the rationale for freezing it, the
vision-language-action lineage the design descends from, and the factored
action space it predicts over. Chapter~\ref{chap:method} details the
problem formulation, the frozen-backbone architecture and trainable action
head, the training objective, the data pipeline, and the
rollout-evaluation protocol. Chapter~\ref{chap:experiments} reports the
$2\times2$ backbone-by-language ablation across both tasks, with
frame-level metrics and in-environment rollout results.
Chapter~\ref{chap:discussion} interprets these results, in particular what
they imply about the contribution of the language pathway and the choice
of frozen backbone. Chapter~\ref{chap:conclusion} summarises the findings
and outlines directions for future work.
